\documentclass[12pt,a4paper]{article}
\usepackage[utf8]{inputenc}
\usepackage[T1]{fontenc}
\usepackage[french]{babel}
\usepackage{amsmath,amssymb}
\usepackage{geometry}
\geometry{margin=2cm}
\usepackage{enumitem}
\usepackage{parskip}
\usepackage{array}
\usepackage{booktabs}
\usepackage{eurosym}   % pour \euro

% --- Cadre de réponse finale ---
\newcommand{\cadreFinal}{%
  \begin{center}
  \fbox{%
    \parbox{0.94\linewidth}{%
      \textbf{Finalement,} \dotfill \\[0.35cm]
      \dotfill
    }%
  }
  \end{center}
}

\begin{document}

\begin{center}
    {\Large \textbf{Fiche élève -- Taux, coefficients multiplicateurs}}\\[0.3cm]
    {\large Barème : Ex1 = 4 pts \quad Ex2 = 6 pts \quad Ex3 = 10 pts \quad Total = 20 pts}\\[0.2cm]
    {\small Pénalités hors barème : jusqu'à \(-2\) points (rédaction, orthographe, grammaire, syntaxe, graphie, présentation).}
\end{center}

\vspace{0.5cm}

\begin{tabular}{p{7cm}p{7cm}}
    Nom : \dotfill & Prénom : \dotfill \\[0.4cm]
    Classe : \dotfill & Date : \dotfill \\
\end{tabular}

\vspace{0.5cm}

\section*{Exercice 1 -- Notions de base (4 points)}

Un commerçant vend un article à \(80\)~\euro{}.

\begin{enumerate}[label=\textbf{\arabic*.}, itemsep=1.2cm]
    \item Le prix augmente de \(15\%\). Calculer le coefficient multiplicateur associé, puis le nouveau prix. \hfill \textbf{(2 pts)}
    
    \vspace{3cm}
    \cadreFinal
    
    \item Le prix de cet article passe ensuite de \(92\)~\euro{} à \(78{,}20\)~\euro{}. Calculer le taux d'évolution en pourcentage. Interpréter le signe. \hfill \textbf{(1 pt)}
    
    \vspace{2.5cm}
    \cadreFinal
    
    \item Un autre article subit une baisse de \(8\%\). Son prix final est \(46\)~\euro{}. Retrouver son prix initial. \hfill \textbf{(1 pt)}
    
    \vspace{2.5cm}
    \cadreFinal
\end{enumerate}

\section*{Exercice 2 -- Évolutions successives et taux global (6 points)}

Une entreprise compte \(15\,000\) clients au 1\textsuperscript{er} janvier 2020.

\begin{itemize}
    \item En 2020, elle perd \(6\%\) de clients.
    \item En 2021, elle gagne \(10\%\).
    \item En 2022, elle perd \(4\%\).
\end{itemize}

\begin{enumerate}[label=\textbf{\arabic*.}, itemsep=1.2cm]
    \item Donner les coefficients multiplicateurs \(CM_1\), \(CM_2\) et \(CM_3\). \hfill \textbf{(1 pt)}
    
    \vspace{2cm}
    \cadreFinal
    
    \item Calculer le coefficient multiplicateur global. \hfill \textbf{(1,5 pt)}
    
    \vspace{2.5cm}
    \cadreFinal
    
    \item En déduire le taux global d'évolution sur les trois années. \hfill \textbf{(1 pt)}
    
    \vspace{2.5cm}
    \cadreFinal
    
    \item Calculer le nombre de clients fin 2022 (arrondir à l'unité). \hfill \textbf{(1,5 pt)}
    
    \vspace{2.5cm}
    \cadreFinal
    
    \item Un stagiaire affirme : « Le taux global est \(-6+10-4 = 0\%\), donc le nombre de clients n'a pas changé. » Expliquer pourquoi il a tort. \hfill \textbf{(1 pt)}
    
    \vspace{3cm}
    \cadreFinal
\end{enumerate}

\section*{Exercice 3 -- Taux moyen et taux inconnu (10 points)}

Un capital de \(2\,000\)~\euro{} est placé. Il subit les évolutions annuelles suivantes :
\[
+10\%,\quad -5\%,\quad +8\%,\quad -3\%,\quad +6\%.
\]

\begin{enumerate}[label=\textbf{\arabic*.}, itemsep=1.2cm]
    \item Calculer le coefficient multiplicateur global sur ces 5 années. \hfill \textbf{(2 pts)}
    
    \vspace{3cm}
    \cadreFinal
    
    \item Calculer le taux global d'évolution en pourcentage (arrondir à \(0{,}01\%\)). \hfill \textbf{(1,5 pt)}
    
    \vspace{2.5cm}
    \cadreFinal
    
    \item Calculer le taux annuel moyen équivalent (arrondir à \(0{,}01\%\)). \hfill \textbf{(3 pts)}
    
    \vspace{3.5cm}
    \cadreFinal
    
    \item Calculer le capital final (arrondir au centime). \hfill \textbf{(1,5 pt)}
    
    \vspace{2.5cm}
    \cadreFinal
    
    \item Un autre placement a augmenté de \(44\%\) en deux ans, avec le même taux annuel \(t\) chaque année. Déterminer \(t\). \hfill \textbf{(2 pts)}
    
    \vspace{3cm}
    \cadreFinal
\end{enumerate}

\newpage

\section*{Grille de correction (réservée au correcteur)}

\begin{center}
\renewcommand{\arraystretch}{1.6}
\begin{tabular}{|>{\raggedright\arraybackslash}p{9cm}|c|c|}
\hline
\textbf{Compétence évaluée} & \textbf{Barème} & \textbf{Points obtenus} \\
\hline
\multicolumn{3}{|l|}{\textbf{Exercice 1 -- Notions de base (4 pts)}} \\
\hline
Calculer un coefficient multiplicateur à partir d'un taux & 1 & \hspace{1.5cm} \\
\hline
Appliquer la relation \(V_f = CM \times V_i\) & 1 & \hspace{1.5cm} \\
\hline
Calculer un taux d'évolution à partir de \(V_i\) et \(V_f\) & 1 & \hspace{1.5cm} \\
\hline
Interpréter un taux d'évolution & 0,5 & \hspace{1.5cm} \\
\hline
Déterminer une valeur initiale connaissant \(V_f\) et \(CM\) & 0,5 & \hspace{1.5cm} \\
\hline
\textbf{Sous-total Exercice 1} & \textbf{4} & \hspace{1.5cm} \\
\hline
\multicolumn{3}{|l|}{\textbf{Exercice 2 -- Évolutions successives et taux global (6 pts)}} \\
\hline
Convertir des taux en coefficients multiplicateurs & 1 & \hspace{1.5cm} \\
\hline
Multiplier des coefficients multiplicateurs pour obtenir un coefficient global & 1,5 & \hspace{1.5cm} \\
\hline
Calculer un taux global à partir du coefficient global & 1 & \hspace{1.5cm} \\
\hline
Calculer une valeur finale après évolutions successives & 1,5 & \hspace{1.5cm} \\
\hline
Argumenter sur la non-additivité des taux & 1 & \hspace{1.5cm} \\
\hline
\textbf{Sous-total Exercice 2} & \textbf{6} & \hspace{1.5cm} \\
\hline
\multicolumn{3}{|l|}{\textbf{Exercice 3 -- Taux moyen et taux inconnu (10 pts)}} \\
\hline
Calculer un coefficient multiplicateur global & 2 & \hspace{1.5cm} \\
\hline
Calculer un taux global en pourcentage & 1,5 & \hspace{1.5cm} \\
\hline
Calculer un taux moyen à partir du coefficient global et du nombre de périodes & 3 & \hspace{1.5cm} \\
\hline
Calculer une valeur finale à partir de \(V_i\) et \(CM_{\text{global}}\) & 1,5 & \hspace{1.5cm} \\
\hline
Résoudre une équation de la forme \(\left(1+\frac{t}{100}\right)^n = CM\) pour trouver un taux inconnu & 2 & \hspace{1.5cm} \\
\hline
\textbf{Sous-total Exercice 3} & \textbf{10} & \hspace{1.5cm} \\
\hline
\multicolumn{2}{|r|}{\textbf{Total barème (sur 20)}} & \hspace{1.5cm} \\
\hline
\end{tabular}
\end{center}

\vspace{0.8cm}

\section*{Pénalités hors barème (jusqu'à \(-2\) points)}

\begin{center}
\renewcommand{\arraystretch}{1.5}
\begin{tabular}{|>{\raggedright\arraybackslash}p{4.2cm}|>{\raggedright\arraybackslash}p{7.5cm}|c|c|}
\hline
\textbf{Critère} & \textbf{Indicateurs} & \textbf{Malus max} & \textbf{Points retirés} \\
\hline
Rédaction & Absence de phrase réponse, pas de justification, calculs non explicités & \(-0{,}5\) & \hspace{1cm} \\
\hline
Orthographe & Fautes d'orthographe d'usage ou grammaticale & \(-0{,}5\) & \hspace{1cm} \\
\hline
Grammaire / Syntaxe & Phrases incorrectes, accords manquants, ponctuation absente & \(-0{,}5\) & \hspace{1cm} \\
\hline
Graphie & Écriture illisible, chiffres ambigus, ratures non assumées & \(-0{,}25\) & \hspace{1cm} \\
\hline
Présentation générale & Copie sale, exercices non identifiés, absence de numérotation & \(-0{,}25\) & \hspace{1cm} \\
\hline
\textbf{Total pénalités} & & \(\mathbf{-2}\) & \hspace{1cm} \\
\hline
\end{tabular}
\end{center}

\vspace{0.5cm}

\begin{center}
\fbox{%
\parbox{0.9\linewidth}{%
\textbf{Calcul de la note finale :}\\[0.2cm]
Note finale \(=\) Total barème (sur 20) \(-\) Total pénalités.\\[0.2cm]
Si le résultat est négatif, la note est ramenée à \(0\).%
}%
}
\end{center}

\vspace{0.5cm}
\noindent\emph{Appréciation :} \dotfill

\end{document}